\documentclass[12pt]{article}

\usepackage[spanish]{babel}
\usepackage[utf8]{inputenc}
\usepackage[T1]{fontenc}
\usepackage{geometry}
\usepackage{setspace}
\usepackage{graphicx}
\usepackage{booktabs}
\usepackage{amsmath}
\usepackage{hyperref}
\usepackage{natbib}

\geometry{margin=1in}
\onehalfspacing

\title{\textbf{Borrador de working paper}\\
Efectos de cercanía a Áreas Naturales Protegidas sobre bienestar local}
\author{Equipo de investigación}
\date{\today}

\begin{document}
\maketitle

\begin{abstract}
% Placeholder: resumir pregunta, datos, estrategia e implicaciones.
\end{abstract}

\section{Introducción}
% Placeholder: motivación, pregunta de investigación y contribución.

\section{Contexto}
% Placeholder: contexto institucional de ANPs y territorio.

\section{Datos}
% Placeholder: fuentes, construcción de muestra, ventana 2014m1--2018m12 y extensión 2019+.

\section{Estrategia empírica}
% Placeholder: definición de tratamiento/control, PSM y especificaciones.

\section{Resultados}
% Placeholder: resultados principales (sin inventar cifras en esta plantilla).

\section{Robustez}
% Placeholder: sensibilidades, extensiones 2019--2022 y pruebas adicionales.

\section{Discusión e implicaciones de política}
% Placeholder: mecanismos, externalidades y relevancia para política pública.

\section{Conclusión}
% Placeholder: hallazgos, límites y agenda de investigación.

\bibliographystyle{apalike}
\bibliography{../bibliography/refs}

\appendix
\section{Apéndice}
% Placeholder: detalles técnicos, tablas y figuras adicionales.

\end{document}
